\documentclass[a4paper, 10pt]{article}

%\usepackage{scalefnt}
%\usepackage{parcolumns}

\usepackage{newclude}
% Zum Einbinden in die Zusammenfassungs-files

\usepackage{amsmath,amsthm,amsfonts,amssymb} % Verbesserter Mathesatz
\usepackage{algorithm2e}
\usepackage{bigfoot} % komplexe Fußnotenapparate(Fußnoten in Fußnoten und andere Späße)%
\usepackage{colortbl}
\usepackage[T1]{fontenc} % normaler erweitere Zeichnesatz
\usepackage{framed, color}  % Ramenpaket für zum Einfügen von schönen Ramen
\usepackage{graphicx}
\usepackage{hyperref} %used for link to creative commons license
\usepackage{listings} %for code-listings (inkl. Tab-Styling)
\usepackage{marvosym}
\usepackage{marginnote}
\usepackage{microtype} % div. Verbesserungen des Schriftsatzes (Grauwert, opt. Randausgleich, Zeilenumbruch)
%\usepackage{multirow}
\usepackage{multicol} %use this with next line for vertical divided environment
%\setlength\columnseprule{.4pt}:465
\usepackage[ngerman]{babel} % Neue Rechtschreibung
\usepackage{pifont}
\usepackage[sans]{dsfont} %für alternative Mengensymbole
\usepackage{stmaryrd} %u.a. für \lightning
\usepackage{tikz} % für Diagramme(Dia!) und Bilder (z.B. *.eps)/für ER-Diagramme/für UML-Diagramme
%\usepackage{tikz-er2} %  er-diagramme
\usepackage{tikz-uml} % uml-Diagramme
\usepackage{units} % z.B. fuer \nicefrac{}{}
\usepackage[utf8]{inputenc} % utf8 für den Editor
\usepackage{wasysym} %u.a. für \lightning
\usepackage{xcolor}


\usetikzlibrary{shapes,decorations,arrows,fit,backgrounds} %Zum diversen zeichen%
%\usetikzlibrary{automata} % für den CFlipper, wenn es den soweit ist			%
\usetikzlibrary{positioning} % positionierung
%\usetikzlibrary{shadows} % fuer schoene schlagschatten
\usetikzlibrary{automata} % für Automaten

%%%%%%%%%%%%%%%%%%%%%%%%%%%
%  Formatierung der Seite
%%%%%%%%%%%%%%%%%%%%%%%%%%%
\usepackage{fancyhdr}
\pagestyle{fancy}		% für den footer										%
\renewcommand{\headrulewidth}{0pt} % damit oben kein dummer Strich kommt		%
\fancyhead{}
\topmargin -2cm 		% Oberer Rand											%
\textheight 25cm		% Texthöhe												%
\textwidth 16.0 cm		% Textbreite											%
\oddsidemargin -0.1cm 	% Warum?												%
\newcommand{\Gruppe}[2]
{
	\lfoot{#1}
	\rfoot{#2}
}
\colorlet{shadecolor}{gray!25} % Farbe für graue Box definieren
%%%%%%%%%%%%%%%%%%%%%%%%%%%%%%%%%%%%%%%%%%%%%%%%%%%%%%%%%%%%%%%%%%%%%%%%%%%%%%%%%
%Farben die definiert werden zum schreiben und zeichnen							%
%%%%%%%%%%%%%%%%%%%%%%%%%%%%%%%%%%%%%%%%%%%%%%%%%%%%%%%%%%%%%%%%%%%%%%%%%%%%%%%%%
\xdefinecolor{schwarz}{HTML}{000000}
\xdefinecolor{dunkelGruen}{HTML}{007D00}
\xdefinecolor{dunkelBlau}{HTML}{0000A0}
\xdefinecolor{dunkelRot}{HTML}{A00000}
\xdefinecolor{dunkelGelb}{HTML}{FFAA00}
\xdefinecolor{hellesGelb}{HTML}{FFCC00}
\colorlet{dGreen}{dunkelGruen}
\colorlet{dBlue}{dunkelBlau}
\colorlet{dRed}{dunkelRot}
\colorlet{dYellow}{dunkelGelb}

%%%%%%%%%%%%%%%%%%%%%%%%%%%%%%%%%%%%%%%%%%%%%%%%%%
%Farbliche Ausgaben:
%Parameter #1: Text oder Mathematische formel...
%z.B. : \gruen{Hallo Welt Test!}
%%%%%%%%%%%%%%%%%%%%%%%%%%%%%%%%%%%%%%%%%%%%%%%%%%

\newcommand{\yellow}[1]{\textcolor{dYellow}{#1}}
\newcommand{\gray}[1]{\textcolor{gray}{#1}}
\newcommand{\red}[1]{\textcolor{red}{#1}}
\newcommand{\green}[1]{\textcolor{green}{#1}}
\newcommand{\blue}[1]{\textcolor{blue}{#1}}
\newcommand{\dGreen}[1]{\textcolor{dGreen}{#1}}
\newcommand{\dBlue}[1]{\textcolor{dBlue}{#1}}
\newcommand{\dRed}[1]{\textcolor{dRed}{#1}}

%%%%%%%%%%%%%%%%%%%%%%%%%%%%%%%%%%%%%%%%%%%%%%%%
%Konfiguration für das darstellen von Quelltext
%%%%%%%%%%%%%%%%%%%%%%%%%%%%%%%%%%%%%%%%%%%%%%%%
\lstset
{
	language=Java, % oder C++, Pascal, {[77]Fortran}, ...
	numbers=left, % Position der Zeilennummerierung
	firstnumber=auto, % Erste Zeilennummer
	basicstyle=\ttfamily, % Textgröße des Standardtexts
	keywordstyle=\ttfamily\color{dRot}, % Formattierung Schlüsselwörter
	commentstyle=\ttfamily\color{dGruen}, % Formattierung Kommentar
	stringstyle=\ttfamily\color{dBlau}, % Formattierung Strings
	numberstyle=\tiny, % Textgröße der Zeilennummern
	stepnumber=1, % Angezeigte Zeilennummern
	numbersep=5pt, % Abstand zw. Zeilennummern und Code
	aboveskip=15pt, % Abstand oberhalb des Codes
	belowskip=11pt, % Abstand unterhalb des Codes
	captionpos=b, % Position der Überschrift
	xleftmargin=10pt, % Linke Einrückung
	frame=single, % Rahmentyp
	breaklines=true, % Umbruch langer Zeilen
	showstringspaces=false, % Spezielles Zeichen für Leerzeichen
	tabsize=2,
	texcl=true
}

%%%%%%%%
% Kopf
%%%%%%%%
\newcommand{\Header}[3]
{
	{\footnotesize \parindent0em
		{\sc Universität Konstanz}                \hfill #1 \\
		{\sc Fachbereich Informatik \& Informationswissenschaft} \hfill #2 \\
		#3 \hfill \today
	}
}

%%%%%%%%%%%%%%%%%%%%%%%
% load some java code
% \loadJava{file}
%%%%%%%%%%%%%%%%%%%%%%%
\newcommand{\loadJava}[1]
{
	\lstinputlisting[language=Java]{#1.java}
}

%%%%%%%%%%%%%%%%%%%%%%%
% load some cpp code
% \loadCpp{file.cpp}
%%%%%%%%%%%%%%%%%%%%%%%
\newcommand{\loadCpp}[1]
{
	\lstinputlisting[language=C++]{#1}
}

%%%%%%%%%%%%%%%%%%%%%%%%%%%%%
% load some code
% \loadCode{Python}{file.py}
%%%%%%%%%%%%%%%%%%%%%%%%%%%%%
\newcommand{\loadCode}[2]
{
	\lstinputlisting[language=#1]{#2}
}

%%%%%%%%%%%%%%%%%%%%%%%%%%%%%%%%%%%%%%%%%%%%%%%%%%%%%%%%%%%%%%%%%%%%%%%
% some symbols
%%%%%%%%%%%%%%%%%%%%%%%%%%%%%%%%%%%%%%%%%%%%%%%%%%%%%%%%%%%%%%%%%%%%%%%
\newcommand{\correct}{\green{\text{\ding{52}}}} %for use in text and math
\newcommand{\wrong}{\red{\text{\ding{56}}}} %for use in text
\newcommand{\tflash}{$\yellow{\lightning}$} %for use in text
\newcommand{\mflash}{\yellow{\lightning}} %for use in math
\newcommand{\follows}{$\Rightarrow$} %used so often...
\newcommand{\good}{\item[\dGreen{\ding{58}}]} %an item with a green plus as bullet point
\newcommand{\bad}{\item[\red{\Emailct}]} %better icon for bad items
\newcommand{\note}[1]{\red{\marginnote{#1}}} %add a red margin note
\newcommand{\fitem}{\item[\follows]} %items with a follows arrow
\newcommand{\hm}{\ensuremath{\overset{-\mkern-11mu-\mkern-3.5mu\rhook}{\smash{\odot}\rule{0ex}{.46ex}}\underline{\hspace{0.5em}}\overset{-\mkern-11mu-\mkern-3.5mu\rhook}{\smash{\odot}\rule{0ex}{.46ex}}}}

%%%%%% make emph bold instead of italic %%%%%
\makeatletter
\DeclareRobustCommand{\em}{%
  \@nomath\em \if b\expandafter\@car\f@series\@nil
  \normalfont \else \bfseries \fi}
\makeatother

%%%%%%%%%%%%%%%%%%%%%%%%%%%%%%%%%%%%%%
% languages for \loadCode
%ABAP		IDL				Plasm
%ACSL		inform			POV
%Ada		Java			Prolog
%Algol		JVMIS			Promela
%Ant		ksh				Python
%Assembler	Lisp			R
%Awk		Logo			Reduce
%bash		make			Rexx
%Basic		Mathematica1	RSL
%C			Matlab			Ruby
%C++		Mercury			S
%Caml		MetaPost		SAS
%Clean		Miranda			Scilab
%Cobol		Mizar			sh
%Comal		ML				SHELXL
%csh		Modula-2		Simula
%Delphi		MuPAD			SQL
%Eiffel		NASTRAN			tcl
%Elan		Oberon-2		TeX
%erlang		OCL				VBScript
%Euphoria	Octave			Verilog
%Fortran	Oz				VHDL
%GCL		Pascal			VRML
%Gnuplot	Perl			XML
%Haskell	PHP				XSLT
%HTML		PL/I
%%%%%%%%%%%%%%%%%%%%%%%%%%%%%%%%%%%%%%

\usetikzlibrary{trees}

\begin{document}

\Gruppe{Stephan Heidinger}{SEng - Zusammenfassung v0.1}
\Header{Software Engineering}{Wintersemester 2011/2012}{Stephan Heidinger}

\begin{shaded}
Dieses Dokument wurde unter der Creative Commons - Namensnennung-NichtKommerziell-Weitergabe unter gleichen Bedingungen (\textbf{CC by-nc-sa}) veröffentlicht. Die Bedingungen finden sich unter \href{http://creativecommons.org/licenses/by-nc-sa/3.0/de}{diesem Link}. \\
\centerline{\includegraphics[scale=1]{../cc-by-nc-sa.png} }
\end{shaded}

\textit{Find any errors? Please send them back, I want to keep them!}

\begin{shaded}
\textbf{Definitions for Software Engineering} \\
``Multi-person construction of multi-version software'' (D. Parnas, 1987)\\
``The application of a systematic, disciplined, quantifiable approach to the development, operation, and maintenance of software; that is, the application of engineering to software.''
\end{shaded}

\section*{Software Crisis}
Term from NATO study group in 1967.
\begin{description}
	\item[fault] mistake/error made by a human during a software activity (erroneous design, requirements, coding)
	\item[failure] observed departure of a system from the desired state
\end{description}
\subsection*{symptoms}
\begin{itemize}
	\item products are delivered late \follows high cost
	\item projects exceed budgeds \follows high cost \follows waste of resources
	\item product doesn't do, whats its supposed to do \follows inefficient \follows hight cost
	\item products are defective \follows high cost (failures, maintenance,\dots), ethical considerations
	\item projects get abandonded before delivery \follows waste of resources
\end{itemize}

\subsection*{Characteristics of software}
\begin{itemize}
	\item Software is \emph{engineered}, not manufactured (custom built, little to no (mechanical) assembly, human intensive process)
	\item Software does not wear out (but: change of requirements, changes in environment)
	\item software is \emph{complex}
	\item software is \emph{detemining system factor} (up to 80\% of development effort)
\end{itemize}

\subsection*{Why is software difficult to produce?}
\begin{itemize}
	\item no similiar system built so far (problems unknown, assumptions about environment may be wrong)
	\item \emph{Requirements} are not well understood/phrased
	\item Requirements \emph{change} during development
	\item complex \emph{interaction}
	\item nature of systems:
		\begin{itemize}
			\item concurrent systems (races, deadlocks, \dots)
			\item embedded systems (hardware interaction, timing, \dots)
			\item information systems (complexity, legacy, \dots)
		\end{itemize}
	\item software is \emph{easily changeble} \follows ``code and fix''
	\item software is \emph{discreet} (either fails, or doesn't)\marginnote{Stimmt auch nicht immer\dots}
\end{itemize}

\subsection*{Software Debvelopment Myths}
\begin{itemize}
	\item Management
		\begin{itemize}
			\item standard books, software, tools, \dots \follows software hard to standardize
			\item behind schedule? add programmers! \follows effort to integrate new people
		\end{itemize}
	\item Customer
		\begin{itemize}
			\item general statement of objectives sufficient
			\item change (as in requirements) can easily be implemented
		\end{itemize}
	\item Pracitioner
		\begin{itemize}
			\item once programm is running, job is done
			\item until programm is running, no way to check quality
			\item only deliverable is working programm
		\end{itemize}
\end{itemize}

\subsection*{Software Engineering}
Communication between a huge number of parties (customer, end user, sw designer, developer, \dots) must be maintained. \\
Highly complex systems: SE has to provide solutions.

\begin{description}
	\item[\follows] base sw production on engineering approach
	\item[\follows] design sw as a \emph{process} that supports
		\begin{itemize}
			\item correctness \& dependability
			\item cost-effectivenes
			\item complexity of project
			\item longevity of product, life cycle, changes
			\item communication amongst parties
		\end{itemize}
	\item[\follows] software engineering process model
\end{description}

\section*{waterfall model}
\begin{tikzpicture}[node distance=50]

\node[draw] (1) at (0,0) {System Design};
\node[draw] (2) [below right of = 1] {Requirements};
\node[draw] (3) [below right of = 2] {Design};
\node[draw] (4) [below right of = 3] {Implementation};
\node[draw] (5) [below right of = 4] {Integration};
\node[draw] (6) [below right of = 5] {Maintenance};

\node (c1) [right of = 1] {};
\node (c2) [right of = 2] {};
\node (c3) [right of = 3] {};
\node (c4) [right of = 4] {};
\node (c5) [right of = 5] {};
\node (c6) [right of = 6] {};

\node (1c) [left of = 1] {};
\node (2c) [left of = 2] {};
\node (3c) [left of = 3] {};
\node (4c) [left of = 4] {};
\node (5c) [left of = 5] {};
\node (6c) [left of = 6] {};

\path[draw,->] (1) .. controls (c1) .. (2);
\path[draw,->] (2) .. controls (c2) .. (3);
\path[draw,->] (3) .. controls (c3) .. (4);
\path[draw,->] (4) .. controls (c4) .. (5);
\path[draw,->] (5) .. controls (c5) .. (6);

\node[node distance=10] (a1) [right of = c1] {};
\node[node distance=10] (a2) [right of = c2] {\textsc{Srs}};
\node[node distance=10] (a3) [right of = c3] {Design};
\node[node distance=10] (a4) [right of = c4] {\textcolor{gray}{Test}};
\node[node distance=10] (a5) [right of = c5] {\textcolor{gray}{Test}};
\node[node distance=10] (a6) [right of = c6] {};

\path[draw,dashed,color=gray,->] (2) .. controls (2c) .. (1);
\path[draw,dashed,color=gray,->] (3) .. controls (3c) .. (2);
\path[draw,dashed,color=gray,->] (4) .. controls (4c) .. (3);
\path[draw,dashed,color=gray,->] (5) .. controls (5c) .. (4);
\path[draw,dashed,color=gray,->] (6) .. controls (6c) .. (5);

\node[node distance=20] (1a) [left of = 1c] {};
\node[node distance=20] (2a) [left of = 2c] {\textcolor{red}{``what''}};
\node[node distance=20] (3a) [left of = 3c] {\textcolor{red}{``how''}};
\node[node distance=20] (4a) [left of = 4c] {};
\node[node distance=20] (5a) [left of = 5c] {};
\node[node distance=20] (6a) [left of = 6c] {};

\end{tikzpicture}

\begin{description}
	\item[System Design] problem \& System requirements, fesability study, define main subsystems (allocate to hw/sw), system design document (\textsc{Sdd}, informal, with customer, somtimes: user manuals, user interfaces, test plans)
	\item[Requirements] requirement analysis (``what'' the system has to do (not how), software requirements describing observable external behaviour (functional, nonfunctional)), software requirement specification (\textsc{Srs})
	\item[Design] ``how'' the system is working, architectural (high level) design (decompose problem into components, global data structures, internal interfaces), detailed design (algorithmic design, internal data structures, programming language(s))
	\item[Implementation] (and Testing) \follows translate design modules into code, test modules in isolation
	\item[Integration] (and Testing) integrate tested modules to form system, \emph{integration testing}, validation, customer accptance tests
	\item[maintenance] deploy product, perform maintenance (corrective, adaptive, perfective), retire product
\end{description}
\subsubsection*{Benefits}
\begin{itemize}
	\item Definition of seperate Tasks (Seperation of concerns)\\dividing complex design problems into smaller units \follows facilitates team work/reproducibility
	\item specification and documentation  \follows enforces documentation, facilitates testing (against requirements specification)
	\item further concepts
		\begin{itemize}
			\item verification / validation (compare intermediate results with requirements and design)
			\item prototyping (mock-up available early, reduces risk)
			\item evolutionary process model (accommodate change)
		\end{itemize}
\end{itemize}

\subsubsection*{Defects}
\begin{itemize}
	\item no feedback loops
	\item document driven, inflexible
	\item large time gap between inception and completion
\end{itemize}

\subsubsection*{Alternatives}
\begin{itemize}
	\item modified waterfall model
	\item V-Model, V-Model XT
	\item Evolutionary Process models
	\item Spiral model
	\item Rational Unified Process
	\item Agile Processes
	\item \dots
\end{itemize}

\subsection*{Software Qualities}
\begin{description}
	\item[Correctness:] the software behaves according to the \emph{requirement specification}
	\item[Reliability:] the software guarantees a level of quality service ($<$Correctness)
	\item[Robustness:] the software behaves ``reasonably'' even in unexprected circumstances (e.g. input, power failure, \dots)
	\item[Maintainability:] software is easy to maintain/extend
	\item[Performance:] system is usable (e.g. input response)
	\item[Reusability:] re-use of previously used, tested and verified code
	\item[Interoperability:] standardisation, interfaces, \dots
\end{description}

\subsection*{Software Engineering Principles}
\begin{description}
	\item[Rigor:] use a method and apply it rigorously to every step
	\item[Formality:] use (mathematical) formalizable methods and notations
	\item[Speration of Concerns:] deal with different aspects in seperate steps
	\item[Abstraction:] seperate the concern of the important aspects from the concern of unimportant ones
	\item[Modularity:] decompose problem into independent modules \follows separation of concern
	\item[Generality:] focus on the discovery of more general problems
\end{description}

\section*{Requiremnt \& Analysis}

\subsection*{Requirement Elicitation}
\begin{shaded}
``A \textbf{requirement} is a condition or capability that must be met or professed by a system component to satisfy a contract, standard, specification or other formally inposed document.'' \\
(ANSI/IEEE Standard 729-1983)
\end{shaded}
requirements \follows description of the \emph{externally visible (interface) \marginnote{not environment behaviour}behaviour}, the ``what''\\
Interaction between the system and the system-relevant portion of the envrionment

\begin{description}
	\item[User (Client) requirements:] statements in natural language plus diagramms
	\item[System requirements:] structured document, detailed description of system's functions, services and operational constraints; may be part of contract between client and contractor
\end{description}
\subsubsection{Types of requirements}
\begin{description}
	\item[functional requirements:] high-level-``what'' the system should do
		\begin{itemize}
			\item statement which services the system should provide
			\item interaction between system and environment (system states, i/o)
			\item how the system should behave in particular situations
			\item describe the system services in detail
			\item often included with ``shall/should''
		\end{itemize}
	\item[nonfunctional requirements:] (if not met, system may be useless (not usable)) \\
	nonfunctional requirements often have quantitative measures (MB, training time, availability, \dots)\\
	\begin{tikzpicture}[
	dist0/.style={},
	dist1/.style={sibling distance=11em},
	dist2s/.style={sibling distance=5.2em},
	dist2b/.style={sibling distance=6.4em},
	down/.style={level distance=7ex},
	downB/.style={level distance=20ex},
	level0/.style={draw,rectangle,fill=blue!20,rounded corners=.8ex},
	level1/.style={draw,rectangle,fill=green!20,rounded corners=.8ex},
	level2/.style={draw,rectangle,fill=yellow!20,rounded corners=.8ex},
	level3/.style={draw,rectangle,fill=gray!20,rounded corners=.8ex},
]

	\coordinate
	child[grow=down,anchor=south, level distance=0ex] {node[level0] {non-functional}}
	[edge from parent fork down]
		child[dist1,down] {node[level1] {product}
			child[dist2s,down] {node[level2] {usability}
			}
			child[dist2s,down] {node[level2] {efficiency}
				child[dist2s,down] {node[level3] {performance}}
				child[dist2s,down] {node[level3] {space}}
			}
			child[dist2s,down] {node[level2] {reliability}
			}
			child[dist2s,down] {node[level2] {portability}
			}
		}
		child[dist1,down] {node[level1] {organisational}
			child[dist2b,downB] {node[level2] {delivery}
			}
			child[dist2b,downB] {node[level2] {implementation}
			}
			child[dist2b,downB] {node[level2] {standards}
			}
		}
		child[dist1,down] {node[level1] {external}
			child[dist2b,down] {node[level2] {interoperability}
			}
			child[dist2s,down] {node[level2] {ethical}
			}
			child[dist2s,down] {node[level2] {legislative}
				child[dist2s,down] {node[level3] {privacy}}
				child[dist2s,down] {node[level3] {safety}}
			}
		}
;
\end{tikzpicture}
		\begin{itemize}
			\item not primarly related to sytem's funcions/services, but to quality and additinal (quantitative) characteristics
			\item constraints on the services/functions:
			\begin{itemize}
				\item reliability (avalability, integrity, security, safety)
				\item accuracy of results
				\item performance / timing
				\item human-computer interface issues
				\item operating and physical constraints
				\item portability and interoperatbility
				\item reliability
				\item response time
				\item storage requirements
				\item standards \dots
				\item also: particular system, programming language or development method
			\end{itemize}
			\item product requirements (product must behave in particular way (execution time, speed, reliability, \dots))
			\item organisational requirements (standards used, impelementation requirements, \dots)
			\item external requirements (interoperatbility/legislative requirements, \dots)
		\end{itemize}
		\item[domain requirements:] (if not met, system may be unworkable in that domain)
			\begin{itemize}
				\item derived from application domain
				\item characteristics/features of domain
				\item constraints on existing requirements
				\item define specific computations
				\item may be expressed in domain specific language \follows often hard to understand
				\item often implicit \follows hard to aquire
			\end{itemize}
\end{description}


\subsubsection*{Requirement Specification Properties}
\begin{description}
	\item[correctness:] facts in requirement specification \follows required properties of the system
	\item[unambiguity:] all specifications have a single interpretation
	\item[completeness:] 
		\begin{enumerate}
			\item every property required of the system is expressed in the specification\\
			{\small does this include things that are \emph{not permitted?}}
			\item the responses of the software system on all types of possible input values are specified
			\item \dots {\small (more possible)}
		\end{enumerate}
	\item[verifiability:] there exists an effective (manual/automatic) process for checking wether a software product satisfies the required properties (formal verification (mathematical proof) or validation (model checking, testing, simulation)), often not possible for every requirement
	\item[consistency:] no two requirements contradict each other: never $A \wedge \neg A$
	\item[traced:] origin of every requirement is clear
	\item[traceable:] requirement specification is edited so it is easy to reference every single requirement (e.g. numbering)
	\item[design independent:] the requirement specification does not require specific software/architecture/algorithms
\end{description}

\subsubsection{Activities during the requirement Stage}
	\begin{description}
		\item[starting point] 
			\begin{itemize}
				\item customer requirements (abstract)
				\item system specification document (hw \& sw)
			\end{itemize}
		\item[Activities (Kotonya+Sommerville)]
			\begin{itemize}
				\item requirement elicitation\marginnote{Elicitation \follows Herausholen} (interviews, scenarios, market observation, \dots) \\
				determine, which of possibly contradictory requirements are important
				\item requirement documentation and specification (comprehensible requirements document)
				\item requirement validation (consistency, completeness, correspondence of documented requirements and abstract customer or user requirements)
			\end{itemize}
			\begin{tikzpicture}[%
	>=stealth,
	node distance=3.5cm,
	on grid,
	auto,
	io/.style={draw, text width=2.5cm, align=center},
	proc/.style={draw, rounded corners=1.2em, text width=2.5cm, align=center},
	dist/.style={node distance = 3cm},
	dick0/.style={draw,line width=1.5pt},
	dick/.style={dick0,->},
	duenn/.style={draw,line width=1.5pt,dashed,->}
]

\node[io] (root) at (0,0) {customer or user requirements (abstract)};

\node[proc] (2) [below left of  = root] {requirements analysis and negotiation};
\node[proc] (1) [left of       = 2   ] {requirements elicitation};
\node[proc] (3) [below right of = root] {requirements documentation and specification};
\node[proc] (4) [right of      = 3   ] {requirements validation};

\node[io,node distance=2cm] (leaf) [below of = 3] {negotiated and validated requirements};

\node[node distance=1.5cm] (phantom1) [below left of =root] {};
\node[node distance=1.5cm] (phantom2) [below right of =root] {};

\path[dick] (root) -| (1);
\path[dick] (1) -- (2);
\path[dick] (2) -- (3);
\path[dick] (3) -- (4);
\path[dick] (3) -- (leaf);
\path[duenn] (root) -| (4);

\path[dick] (phantom2) -| (1.north east);
\path[dick] (phantom2) -| (2);
\path[dick] (phantom1) -| (3);
\path[dick0] (phantom1) -| (4.north west);


%	\node[state] (A)              {A};
%	\node        (B) [right of=A,fill=blue!25,text width=3cm]{This is a demonstration text for showing how line breaking works.};;
%	\path[->] (A) edge (B);


\end{tikzpicture}
			\item[social activity:]
			Not a single person knows everything about the system \follows communication is needed \follows difficult (technical language, ambiguities, customers not aware of needs, personalities.)
	\end{description}

\subsubsection*{Interviewing / Structured Questioning Techniques}
\begin{enumerate}
	\item context-of-system (why are we building this system/who are the users/critical functionality/needs)
	\item open-ended questions (produce large amount of info/use, when not much known yet
	\item close-ended questions (specific questions))
	\item rephrase questions (make sure answers understood right/inconsistencies/ambiguities)
	\item whom else to interview
\end{enumerate}

\subsubsection*{Learning from Existing Systems}
\begin{itemize}
	\item analyse user manuals
	\item use/play with existing systems
	\item market analysis (competing systems, market research, which features to be included
	\item reverse engineering)
\end{itemize}

\subsubsection*{requirement uncertainty}
When the customer does not know, what he really needs/wants \follows prototype.
Built prototype early (only if prototype build is faster than system build),
use prototype for requirement elicitation.
Prototyping also part of some life cycle and process models (spiral model)

\subsubsection{Requirements Gathering - Brainstorming}
\begin{itemize}
	\item generate ideas (free of criticism/judgement) \follows as many ideas as possible
	\item discuss, revise, organize ideas
	\item evaluate, prioritize
\end{itemize}

\subsubsection{\textsc{Fast} (Facilitated Application Specification Technique)}
\begin{itemize}
	\item overcome we/then attitude (developer/user/customer)
	\item team oriented
	\item guidelines:
		\begin{itemize}
			\item must participate entire meeting
			\item participants are queal
			\item preparation as important as meeting
			\item pre-meeting documents are only ``proposed''
			\item off-site location preferred
			\item set agenda, maintain ist
			\item no technical details
		\end{itemize}
\end{itemize}

\subsubsection{\textsc{Jad} (Joint Application Design, 1977)}
\begin{itemize}
	\item technique to get all participants to aggree on sw requirements and design
	\item 20\% reduction of life cylce cost, \~15\% functionality less missed
	\item \textsc{Jad}/Plan \follows software requirement elicitation
	\item \textsc[Jad]/Design \follows software design
	\item Main Objective:
		\begin{enumerate}
			\item group dynamics \follows enhace capcabilites ofindividuals
			\item communication and understanding \follows use of visual 	aids
			\item rational, organized process \follows repeatability
			\item standardized documentations \follows use of standard forms
		\end{enumerate}
	\item customisation (prepare tasks for session)
	\item session (customized, facilitated meetings (developers and user))
	\item wrap-up (finalize results from sesion)
	\item participants
		\begin{description}
			\item[session leader:] manages sessions, management/communications skills, problem domain competence
			\item[analyst:] \follows session output, technical understanding req.
			\item[executive sponsor:] high-level strategic insight into system, executive-level decissions (resources, \dots)
			\item[user representative:] ``user''
			\item[information systems representative:] feasibility assessment
			\item[specialist:] knowledge from application domain
		\end{description}
	\item Customization phase
	\begin{itemize}
		\item orientation (sponsor endorses project, session leaders (+analyst) gains understanding of system/environment)
		\item organize team (select participants, think of questions ahead)
		\item tailor the process (how much time/resources, customize documents)
		\item prepare materials (slides, presentations, \dots, white boards, flip chats, \dots, agenda)
	\end{itemize}
	\item Session phase
	\begin{itemize}
		\item orientation (welcome, overview)
		\item define high-level requirements (objectives, benefits, strategic/future considerations, constraints and assumptions, security/audit/control) \follows recorded by analyst, discussion (refine, assess requirements)
		\item defince scope of system (organize and prioritize requirements, scopes, s.t. system meets objectives, but not too costly or complex)
		\item prepare \textsc{Jad}/Design (estimate resources, identify participants, schedule meetings)
		\item document issues and considerations (problems that affect requirements \follows assign a person for resolution)
		\item conclusion (review information and decissions, voice remaining concerns, conclude s.t. particip. have sense of ownership/commitment)
	\end{itemize}
	\item Warp-Up Phase
	\begin{itemize}
		\item transform session documents into formal planning documents (analyst)
		\item edit \textsc{Jad}/Plan
		\item review \textsc{Jad}/Plan (all participants \follows changes possible
		\item sponsor approval)
	\end{itemize}
\end{itemize}

\subsection*{Object Oriented Analysis}
mapping onto waterfall process
\begin{itemize}
	\item requirements \follows object-oriented \emph{analysis}
	\item architectural design \follows system design
	\item detailed design \follows object design
\end{itemize}
oo swe processes are \emph{more continuous}
\begin{itemize}
	\item which are \emph{objects of interest}
	\item what are they doing
	\item \textsc{Uml} use case, sequence, class diagrams
\end{itemize}

\begin{shaded}
\textbf{Object} \marginnote{J. Rumbaugh et al}\\
``A discrete entity with a well defined boundary and identity that encapsulates state and behaviour; ...'' \\
``A concept, abstraction or thing with crisp boundaries and meaning for the problem at hand...''\\
``Objects serve two purposes: The promote understanding of the real world and provide a practical basis for computer implementation'' \\
``All objects have identity and are distinguishable.''
\end{shaded}

\subsubsection{Obejct-Oriented Software Engineering}
\begin{itemize}
	\item software is defined with
	\begin{description}
		\item[objects:] entities, correspondence in real world (identity, attributes, methods)
		\item[classes:] abstact of objects
	\end{description}
	\item objects encapsulate data (interface to data)
	\item interface is relevant, not implementation
	\item reuse through inheritance
	\item abstraction through polymorphism
\end{itemize}

\subsubsection{Software Models}
\begin{itemize}
	\item software systems complex \follows break down into subproblems
	\item software models \follows abstraction of real world
	\item usage
		\begin{description}
			\item[eary life-cycle stages:] evaluate properties of real-world system (elicitate, document, verify, validate requirements, simulation)
			\item[design stage:] document architecture, assess performance/behaviour
			\item[implementation/coding stage:] automatically synthesize code (class skeletons, state machine behaviour)
			\item[testing stage] what to test (requirements models)
			\item[maintenance stage:] document existing system for evolution
		\end{description}
\end{itemize}

\subsection*{Analysis with the UML}
\begin{itemize}
	\item visual notation for analysis and design
	\item standardized by \emph{Object Management Group} \url{http://www.omg.org}
\end{itemize}

\subsubsection{Static}
\begin{description}
	\item[Use Case Diagrams]
	\begin{itemize}
		\item often first step i OO-dev-process
		\item a use case represents a usage scenario (atomic)
		\item documents the functinoalities the system provides to user (``what'' does the system do)
		\item actors (anything or anyone interacting with the system, provide input/output) and principal use cases (interaction between actor and system)
		\item \begin{tikzpicture}[
]

\begin{umlsystem}[x=4, fill=red!10]{The system}
\umlusecase{use case1}
\umlusecase[y=-2]{use case2}
\umlusecase[y=-4]{use case3}
\umlusecase[x=4, y=-2,]{use case4}
\umlusecase[x=6, fill=green!20]{use case5}
\umlusecase[x=6, y=-4]{use case6}
\end{umlsystem} 

\umlactor{user}
\umlactor[y=-3]{subuser}
\umlactor[x=14, y=-1.5]{admin}

\umlinherit{subuser}{user}
\umlassoc{user}{usecase-1}
\umlassoc{subuser}{usecase-2}
\umlassoc{subuser}{usecase-3}
\umlassoc{admin}{usecase-5}
\umlassoc{admin}{usecase-6}
\umlinherit{usecase-2}{usecase-1}
\umlHVextend[name=extd]{usecase-4}{usecase-5}
\umlinclude[name=incl]{usecase-3}{usecase-4}

\umlnote[x=7, y=-6]{incl-1}{use case3 uses something of use case4}
\umlnote[x=7,y=2]{extd-1}{use case5 might lead to use case4}

\end{tikzpicture}
		\item <<include>> \follows a use case uses functionality from another use case
		\item <<extend>> \follows exceptional or optional behaviour
	\end{itemize}
	\item[Class Diagrams:]
	\begin{itemize}
		\item Phenomen (object in the world of a domain as perceived
		\item Concept (properties iof a phenomen) is a 3-tuple: name (distringuishes it from others), purpose (properties that determine, if phenomen is member of concept), members (phenomena, that are part of concept)
		\item Type (abstraction in context of programming (name: int) \\
		\texttt{SimpleWatch} or \texttt{Firefighter}
		\item Class (abstraction in the context of oo-languages, encapsulates state(variables) and behaviour(methods))
		\item Instance (existing instance of a class) \underline{\texttt{myWatch:SimpleWatch}} or \underline{\texttt{Joe:Firefighter}}
		\item \begin{tikzpicture}[]

\begin{umlpackage}{package}
\umlclass[template=T]{Template Class}{
  n : uint \\ t : float
}{}
\umlclass[x=4]{real class}{
  d : double
  }{
  \umlvirt{setB(b : B) : void} \\ getB() : B}
\end{umlpackage}


\end{tikzpicture}
		\item Generality (most general aspect of problem)
		\item Abstraction (not individual objects, what they have in common)
		\item 		``The descriptor for a set of objects that share the same attributes, relationships and behavior. A class represents a concept within the system being modeled.''
		\item used during requirement analysis (model problem domain concepts), system design (model subsystems and interfaces), object design (model prog language classes)
		\item Association: ``A semantic condition or restriction represented as an expression.'' \\
		\begin{tikzpicture}[]
\umlclass[x=-3,y=0]{Person}{name: String \\ age: Integer}{}
\umlclass[x=3,y=0]{Company}{name: String}{}
\umlassoc[mult1=0..*,mult2=*,name=contract]{Person}{Company}
\umlclass[x=0,y=2]{WorkContract}{salary: Integer}{terminate()}
\umlassoc[style=dashed]{WorkContract}{contract-1}

\umlclass[x=-3,y=-2.5]{Person}{name: String \\ age: Integer}{}
\umlclass[x=3,y=-2.5]{Car}{make: carMakers \\ color: carColors}{}
\umlassoc[mult1=1,mult2=0..*]{Person}{Car}

\umlclass[x=-3,y=-4.5]{Triangle}{}{}
\umlclass[x=3,y=-4.5]{Corners}{}{}
\umlassoc[mult1=1,mult2=3]{Triangle}{Corners}


\end{tikzpicture}
		\item Dependencies: \\
		\begin{tikzpicture}[]

\umlclass[x=-3,y=0]{BoxOffice}{}{}
\umlclass[x=3,y=0]{SchedulingEngine}{}{}
\umlimpl[stereo=use]{BoxOffice}{SchedulingEngine}

\end{tikzpicture}
		\item Aggregation: \\
		\begin{tikzpicture}[]
\umlclass[x=-0,y=0]{Document}{}{}
\umlclass[x=0,y=-3,name=Toc]{Table of Contents}{}{}
\umlclass[x=3,y=-3]{Section}{}{}
\umlclass[x=-3,y=-3]{Index}{}{}
\umlVHVaggreg[mult2=1,pos=2.5]{Document}{Index}
\umlaggreg[mult2=1..*]{Document}{Toc}
\umlVHVaggreg[mult2=0..1,pos=2.5]{Document}{Section}

\end{tikzpicture}
		\item Composition \\
		\begin{tikzpicture}[]
\umlclass[x=-0,y=0]{Company}{}{}

\umlclass[x=2,y=-3]{Departement}{}{}
\umlclass[x=-2,y=-3]{Office}{}{}
\umlVHVcompo[mult2=1..*,mult1=1,pos=2.5]{Company}{Departement}
\umlVHVcompo[mult2=1..*,pos=2.5]{Company}{Office}
\umlcompo[mult=*, angle1=0, angle2=-90, loopsize=2.5cm]{Departement}{Departement}

\end{tikzpicture}
		\item Generalisation \\
		\begin{tikzpicture}[]
\umlclass[x=-0,y=0]{Staff}{name: String\\SIN: Integer}{calc\_paycheck()}
\umlclass[x=2,y=-4]{SupportStaff}{hours: Integer \\ overtime: Integer \\ hourly\_rate: Integer}{calc\_paycheck()}
\umlclass[x=-2,y=-4]{Engineer}{clearance: Boolean\\salary: Integer}{calc\_paycheck()}
\umlVHVinherit[]{SupportStaff}{Staff}
\umlVHVinherit[]{Engineer}{Staff}

\end{tikzpicture}
		\item Visibility (+ public, - private, \# protected)
	\end{itemize}
\end{description}

\subsubsection{Dynamic}
\begin{description}
	\item[Sequence Diagrams] used during requirement analysis (refinde dynamic behaviour), design (document subsystem interfaces), good for real-time specifications\\
	\begin{tikzpicture}
\begin{umlseqdiag}
%\umlobject[class=simpleWatch,stereo=actor]{myWatch}
%\umlactor[class=simpleWatch]{myWatch}
%\umlentity[class=simpleWatch]{myWatch}
%\umlboundary[class=simpleWatch]{myWatch}
%\umlcontrol[class=simpleWatch]{myWatch}
%\umldatabase[class=simpleWatch]{myWatch}
%\umlmulti[class=simpleWatch]{myWatch}
%\umlobject[class=simpleWatch]{b}

\umlactor[class=WatchUser]{a}
\umlobject[class=SimpleWatch]{b}
\umlobject[class=SomethingElse]{c}

\begin{umlcall}[op={pressButton1()},type=synchron, return=0]{a}{b}

\end{umlcall}

\begin{umlcall}[op={doSomething()}]{a}{c}

\end{umlcall}

\end{umlseqdiag}
\end{tikzpicture}
	\begin{itemize}
		\item classes \follows columns
		\item messages \follows arrows
		\item return values  \follows dashed arrows
		\item activations \follows narrow rectangles
		\item lifelines \follows dashed lines
	\end{itemize}
	\item[State Chart Diagram] state machine description (only for dynamically interesting objects)
	\item[Dynamic Modeling:] Collection of State Chart Diagrams
\end{description}

Example Use Case Format \\
\begin{tabular}{|l|l|}
\hline %%%%%%%%%%%%%%%%%%%%%%%%%%%%%%%%%%%%%%%%%%%%%%%%%%%%%%%%%%%%%
\textbf{use case name}	& ReportEmergency \\
\hline %%%%%%%%%%%%%%%%%%%%%%%%%%%%%%%%%%%%%%%%%%%%%%%%%%%%%%%%%%%%%
\textbf{entry condition}& FieldOfficer activates ``ReportEmergency'' function \\
\hline %%%%%%%%%%%%%%%%%%%%%%%%%%%%%%%%%%%%%%%%%%%%%%%%%%%%%%%%%%%%%
\textbf{flow of events}	& FRIEND responds by presenting form \\
						& FielOfficer fills out form \\
						& Dispatcher reviews information \\
\hline %%%%%%%%%%%%%%%%%%%%%%%%%%%%%%%%%%%%%%%%%%%%%%%%%%%%%%%%%%%%%
\textbf{Exit condition}	& FieldOfficer receives ack and response \\
\hline %%%%%%%%%%%%%%%%%%%%%%%%%%%%%%%%%%%%%%%%%%%%%%%%%%%%%%%%%%%%%
\end{tabular}

\subsubsection{Object Oriented Analysis}
\begin{itemize}
	\item use-case driven approach (followed by most UML-based approaches)
	\item linguistic analysis method (requirements in natural language  \follows cast into OO analysis model, more traditional OO approach
	\item identify Entity Objects (persistent information tracked by system)
	\item indentify Boundary Objects (interactions betwen actor/system)
	\item identify Control Objects (responsible for realizinh use cases)
	\item identify Associations, Aggregates, Attributes
	\item Model State-Independent Behaviour of objects
	\item Mapp Use Cases with Sequence diagrams
	\item Model Inheritance Relationships
	\item Review Analysios Model
\end{itemize}

\subsection*{Document Requirements / SRS}
\begin{shaded}
\textbf{Requirement Specification:} ``A specification that sets forth the
requirements for a system or system component; \dots typically included
are functional requirements, performance requirements, interface
requirements, design requirements and development standards.'' \\
``A \emph{software requirements specification} is a document containing acomplete specification of \textbf{what} the system will do without saying how it will do that.''
\end{shaded}
\begin{itemize}
	\item contractual agreement (\follows litigation)
	\item shall \follows mandatory
	\item should \follows desired
	\item ambiguity \follows to be prevented
	\item IEEE 830-1998
	\item DIN 69905 Lastenheft/Pflichtenheft
	\item SRS by Sommerville \\
	\begin{tabular}{|p{3cm}|p{8cm}|}
	\hline %%%%%%%%%%%%%%%%%%%%%%%%%%%%%%%%%%%%%%%%%%%
	\textbf{Chapter} & \textbf{Description} \\
	\hline %%%%%%%%%%%%%%%%%%%%%%%%%%%%%%%%%%%%%%%%%%%
	Preface & expected readership, version history, rationale for creation of new versions, summary of changes \\
	\hline %%%%%%%%%%%%%%%%%%%%%%%%%%%%%%%%%%%%%%%%%%%
	Introduction & need for the system, systems functions, how system works with other systems, how system fits into buisness \\
	\hline %%%%%%%%%%%%%%%%%%%%%%%%%%%%%%%%%%%%%%%%%%%
	Glossary & define technical terms, assume no knowledge \\
	\hline %%%%%%%%%%%%%%%%%%%%%%%%%%%%%%%%%%%%%%%%%%%
	User requirements definition & services provided for user, nonfunctional also described (natural language, diagrams, understandable notations, process and product standards \\
	\hline %%%%%%%%%%%%%%%%%%%%%%%%%%%%%%%%%%%%%%%%%%%
	System architecture & high-level overview of anticipated system architecture, distribution of functions across system modules, reused components highlighted \\
	\hline %%%%%%%%%%%%%%%%%%%%%%%%%%%%%%%%%%%%%%%%%%%
	System requirements specification & functional \& nonfunctional requirements in detail, define interfaces \\
	\hline %%%%%%%%%%%%%%%%%%%%%%%%%%%%%%%%%%%%%%%%%%%
	System models & graphical system models, relationships between system components, systen and environment, data-flow models, semantic data models \\
	\hline %%%%%%%%%%%%%%%%%%%%%%%%%%%%%%%%%%%%%%%%%%%
	System evolution & fundamental assumptions on which system is based, anticipated changes due to hw evolution, changing user needs, \dots \\
	\hline %%%%%%%%%%%%%%%%%%%%%%%%%%%%%%%%%%%%%%%%%%%
	Appendices & detailed, specific information related to application, hw/database descriptions \\
	\hline %%%%%%%%%%%%%%%%%%%%%%%%%%%%%%%%%%%%%%%%%%%
	Index & alphabetic index, index of diagrams, index of functions, index of tables, \dots \\
	\hline %%%%%%%%%%%%%%%%%%%%%%%%%%%%%%%%%%%%%%%%%%%
	\end{tabular}
	\item IEEE 830-1998 SRS
	\begin{enumerate}
		\item Introduction
		\begin{enumerate}
			\item Purpose
			\item scope
			\item definitions, acronyms, abbreviations
			\item references
			\item overview
		\end{enumerate}
		\item overall description
		\begin{enumerate}
			\item product perspective
			\item product functions
			\item user characteristics
			\item constraints
			\item assumptions and dependencies
			\item apportioning of requirements
		\end{enumerate}
		\item specific requirements
		\begin{enumerate}
			\item external interfaces
			\item functions
			\item performance requirements
			\item logical database Requirements
			\item design constraints
			\item software system attributes
			\item organizing the specific requirements
			\item additional comments
		\end{enumerate}
	\end{enumerate}
	\item Lastenheft (DIN 69905), vom Auftraggeber erstellt
	\begin{enumerate}
		\item Ausgangssituation/Zielsetzung
		\item Produkteinsatz
		\item Produktübersicht
		\item Funtionale Anforderungen
		\item Nichtfunktionale Anforderungen \\
			Benutzbarkeit, Zuverlässigkeit, Effizienz, Änderbarkeit, Übertragbarkeit
		\item Risikoakzeptanz
		\item Skizze des Entwicklungszyklus und der Systemarchitektur
		\item Lieferumfang
		\item Abnahmekriterien
		%2-201
	\end{enumerate}
	\item Pflichtenheft (DIN 69905) \\
	vertraglich bindend, präzise, detailiert, vollständige Beschreibung der zu erfülenden Leistung, Akzeptanztests\\
	vom Aufgragnehmer erstellt, vom Auftraggeber bestätigt
	\begin{enumerate}
		\item Zielbestimmung (Musskriterien, Wunschkriterien, Abgrenzungskriterien (sollen nicht erreicht werden))
		\item Produkteinsatz (Anwendungsbereiche, Zielgruppen, Betriebsbedingungen)
		\item Produktübersicht
		\item Produktfunktionen (genau, detailiert)
		\item Produktdaten (langfristig zu speichernde Daten aus Benutzersicht)
		\item Produktleistungen (Anforderungen Zeit, Genauigkeit)
		\item Qualitätsanforderungen
		\item Benutzeroberfläche (grundlegende Anforderungen, Zugriffsrechte)
		\item Nichtfunktionale Anforderungen (Gesetze, Normen, Sicherheitsanforderungen, Plattformabhängigkeiten)
		\item technische Produktumgebung (SW und HW, organisatorische Rahmenbedingungen, Schnittstellen)
	\end{enumerate}
\end{itemize}

\section*{Design}

\subsection*{Design Objectives}
What \follows How \\
Manage Complexity \\
Design for Change (anticipate change at all stages, reduce cost of changes)

\begin{itemize}
	\item design process should not suffer from ``tunnel vision''
	\item design should be traceable to analysis model
	\item don't reinvent the wheel
	\item minimize intelectual distance between software and real-world problem
	\item uniformity and integration
	\item acommodate change
	\item degrade gently, when aberrant data/events/\dots encountered
	\item Design is not coding
	\item assess design for quality while it is created
	\item review design
\end{itemize}

\subsection*{Classical and Object-Oriented Design}
\subsubsection*{Procedural Design}
\begin{itemize}
	\item decomposition of system into modules
	\item describe module interactions
	\item structure of procedure definitions/invokations
	\item no abstract data types/inheritance
\end{itemize}
\begin{description}
	\item[Architectural Design:] determine major elemts (modules) and relationships
	\item[Interface Design:] how to design interrelation/communiccation/mutural reliance (contracts) of modules
	\item[Data Design:] data structures
	\item[Procedural Design:] determine procedural description for elements
\end{description}

\paragraph*{Module}
\begin{itemize}
	\item well defined component of system (provides set of services to other modules, consists of (name, interface, body), information hiding)
	\item master complexity, facilitate documentation, enable teamwork
	\item interface descrition (``contract server/client'') languages (CORBA IDL, Eiffel)
	\item \emph{Cohesion} (intra module bindings, ``degree to which a module performs one and only one function'')
	\item \emph{coupling} (inter module bindings, ``degree to which a module is connected to other modules'' (direct access, data transfer, \dots))
	\item desired: hoch cohesion, low coupling
\end{itemize}

\subsubsection*{Object-Oriented Design}
\begin{itemize}
	\item model system als collection of classes + objects
	\item description of object interfaces
	\item encapsulation, abstract data types
	\item reuse through inheritance
	\item polymorphism
\end{itemize}
\paragraph*{Activities}
\begin{itemize}
	\item partition model into subsystems
	\item identify concurrency
	\item allocate subsystems to processor tasks
	\item develop ui design
	\item choose strategy for data management
	\item identify global resources (and control mechaninsm, task management)
	\item consider boundary conditions (and how to be handled)
	\item review, consider trade-offs
\end{itemize}
\begin{description}
	\item[Booch Method:] macro development process (architectural planning, patritioning, layering) + micrp development process (rules for implementing specifics)
	\item[Rumbaugh Method:] system design (analysis models represent system, layout of components, partitioning into subsystems, taking exec environment into account) + object design (design of algos/data structures)
	\item[Jacboson Method:] traceability from analysis model, adaptrion of analysis model to real-world env, primary design objects categorized (interface, entity, control), determine communication \follows organize into subsystems
	\item[Wirfs-Brock Method:] analysis \follows design, def of protocols through inter-object contracts, specification of classes and subsystems
\end{description}
\begin{itemize}
	\item maintenance (objects work standalone), reuse (objects can be used again), reduce semantic gaps (real-world $\leftrightarrow$ software)
\end{itemize}

\paragraph*{Generalisation \& Inheritance}
\begin{itemize}
	\item objects are members of classes (define attributes, operations)
	\item class hirarchy \follows Generalisation
	\item generalistaion in UML \follows inheritance in OO
	\item abstraction \follows classify entities
	\item inheritance graph \follows organistational knowledge about domain/system
	\item \emph{Cohesion} how well fits the functionality of a class together
	\item \emph{Coupling} B subclass of A \follows B inherits from A \follows coupling between A/B, reflects real-world structure (not design decision)
\end{itemize}

\subsection*{Design with the UML}
\subsubsection*{types of models}
\begin{description}
	\item[design models:] show objects, object classes, relationships
	\item[static models:] describe static structure, object classes and relationships
	\item[dynamic models:] describe dynamic interactions between objects
\end{description}

\subsubsection*{Typically used models}
\begin{description}
	\item[subsystem/component diagrams:] structural refinement of objects \\
	A \textbf{UML component} ``represents a modular part of a system, that encapsulates its content and whose manifestation is replaceable within its environment. A component defines its behavior in terms of provided and required interfaces
'' \follows subsystem \\
	\begin{tikzpicture}


\begin{umlcomponent}{System1}
\umlbasiccomponent{SubSystem11}
\umlbasiccomponent[x=5]{SubSystem12}

\umlassoc{SubSystem11}{SubSystem12}
\end{umlcomponent}

\begin{umlcomponent}[y=-4]{System2}
\umlbasiccomponent{SubSystem21}

\end{umlcomponent}

\umlHVassemblyconnector[interface=12]{SubSystem21}{SubSystem12}

%\umlbasiccomponent[y=15]{B}
%\umlbasiccomponent[y=10]{C}

%\umlprovidedinterface[interface=C-interface]{C}
%\umlrequiredinterface[interface=B-interface,with port, distance=3cm,padding=2.5cm]{B}

%\umlassemblyconnector[interface=AE,with port]{B}{C}

\end{tikzpicture}
	\item[sequence diagrams:] model order of interactions among objects\\
	\begin{tikzpicture}
\begin{umlseqdiag}
%\umlobject[class=simpleWatch,stereo=actor]{myWatch}
%\umlactor[class=simpleWatch]{myWatch}
%\umlentity[class=simpleWatch]{myWatch}
%\umlboundary[class=simpleWatch]{myWatch}
%\umlcontrol[class=simpleWatch]{myWatch}
%\umldatabase[class=simpleWatch]{myWatch}
%\umlmulti[class=simpleWatch]{myWatch}
%\umlobject[class=simpleWatch]{b}

\umlactor[class=WatchUser]{a}
\umlobject[class=SimpleWatch]{b}
\umlobject[class=SomethingElse]{c}

\begin{umlcall}[op={pressButton1()},type=synchron, return=0]{a}{b}

\end{umlcall}

\begin{umlcall}[op={doSomething()}]{a}{c}

\end{umlcall}

\end{umlseqdiag}
\end{tikzpicture}
	\item[statechart diagrams:] model state behaviour of object
	\item[acitivity diagram:] model concurrent behabiour of object
\end{description}

\begin{description}
	\item[use-case diagram:] model interactions (include description)
\end{description}

\subsection*{Interfaces and Contracts}

\subsection*{Software Architecture and Patterns}

\subsection*{From Design to Code}

\subsection*{Documenting Design}

\end{document}
